\documentclass[10pt]{article}
\usepackage[T1]{fontenc}
\usepackage[utf8]{inputenc}
\usepackage[french]{babel}
\usepackage{amsmath,amssymb,mathrsfs,mathpazo}
\usepackage{geometry}
\geometry{a4paper,margin=2cm}
\usepackage{graphicx,array,tabularx,tabularray,colortbl,multirow,multicol,booktabs,enumitem,pifont,ulem,fancybox,diagbox,cancel}
\usepackage{pstricks,pstricks-add,pst-plot,pst-tree,pst-node,pst-text,pst-func,pst-math,pst-bspline,pst-3dplot,pst-eucl,pst-solides3d}
\let\pscalculate\undefined
\usepackage{pst-bezier}
\usepackage{tikz,pgfplots,tkz-tab}
\usetikzlibrary{arrows,arrows.meta,patterns,calc,positioning,trees,decorations.markings,shapes.geometric,shapes.misc,shapes.arrows,intersections,decorations.pathmorphing,decorations.pathreplacing,backgrounds,shadows,babel}
\pgfplotsset{compat=1.18}
\usepgfplotslibrary{fillbetween}
\usepackage[np]{numprint}
\usepackage{listings,makecell,pas-tableur,twoopt}
\let\barre\undefined
\usepackage[pstricks]{tablvar}
\definecolor{darkgreen}{RGB}{0,100,0}
\providecommand{\red}{\color{red}}
\providecommand{\blue}{\color{blue}}
\providecommand{\green}{\color{green}}
\providecommand{\black}{\color{black}}
\providecommand{\white}{\color{white}}
\providecommand{\gray}{\color{gray}}
\providecommand{\cyan}{\color{cyan}}
\providecommand{\magenta}{\color{magenta}}
\providecommand{\yellow}{\color{yellow}}
\providecommand{\violet}{\color{violet}}
\providecommand{\orange}{\color{orange}}
\providecommand{\euro}{\texteuro}
\providecommand{\EUR}{\texteuro}
\providecommand{\ieme}{\textsuperscript{e}}
\providecommand{\degres}{\ensuremath{^\circ}}
\providecommand{\degree}{\ensuremath{^\circ}}
\providecommand{\vv}[1]{\overrightarrow{#1}}
\providecommand{\Coord}[1]{#1}
\providecommand{\decosix}{}
\providecommand{\starredbullet}{\ensuremath{\bullet}}
\pagestyle{empty}
\setlength{\parindent}{0pt}
\usepackage{adjustbox}\DeclareUnicodeCharacter{2236}{:}\DeclareUnicodeCharacter{2190}{\ensuremath{\leftarrow}}
\providecommand{\R}{\mathbb{R}}
\providecommand{\N}{\mathbb{N}}
\providecommand{\Z}{\mathbb{Z}}
\providecommand{\Q}{\mathbb{Q}}
\providecommand{\D}{\mathbb{D}}
\providecommand{\C}{\mathbb{C}}
\providecommand{\e}{{\mathrm{e}}}
\providecommand{\vect}[1]{\overrightarrow{#1}}
\providecommand{\vectt}[1]{\overrightarrow{\mathrm{#1}}}
\providecommand{\vvt}[1]{\overrightarrow{\mathrm{#1}}}
\providecommand{\np}[2][]{{#2~#1}}
\providecommand{\vv}[1]{{\overrightarrow{#1}}}
\renewcommand{\d}{{\mathrm{d}}}
\providecommand{\degres}{{^\circ}}
\providecommand{\degre}{{^\circ}}
\providecommand{\blue}{\color{blue}}
\providecommand{\red}{\color{red}}
\providecommand{\enskip}{\quad}
\providecommand{\llbracket}{\lbrack\!\lbrack}
\providecommand{\rrbracket}{\rbrack\!\rbrack}
\providecommand{\euro}{{\text{€}}}
\providecommand{\no}{{\text{n°}}}
\providecommand{\interval}[2]{{[#1;#2]}}
\providecommand{\intervFO}[2]{{[#1;#2[}}
\providecommand{\dotfill}{\ldots}

\providecommand{\ds}{\displaystyle}
\providecommand{\textding}[1]{#1}
\providecommand{\Oij}{\(\mathrm{O};\vec{\imath},\vec{\jmath}\)}
\providecommand{\Oijk}{\(\mathrm{O};\vec{\imath},\vec{\jmath},\vec{k}\)}
\providecommand{\Ouv}{\(\mathrm{O};\vec{u},\vec{v}\)}
\providecommand{\intervFF}[2]{\ensuremath{[#1;#2]}}
\renewcommand{\euro}{\ifmmode\text{\texteuro}\else\texteuro\fi}
\DeclareUnicodeCharacter{00B0}{\ifmmode^{\circ}\else\textdegree\fi}
\DeclareUnicodeCharacter{20AC}{\ifmmode\text{\texteuro}\else\texteuro\fi}
% Énoncés : Éducation nationale. Transcription et corrigés : APMEP, Denis Vergès et contributeurs. Découpage et classement : CoopMaths, Sébastien Lozano, Éric Elter et contributeurs. Adaptation de la présentation : Maths Explorer.
% CC BY-SA : https://coopmaths.fr/www/about/
% https://coopmaths.fr/alea/static/bac/2024/tex/bac_2024_06_sujet1_metropole_devoile_1.tex
\begin{document}
\noindent\textit{APMEP / CoopMaths. CC BY-SA. Présentation adaptée par Maths Explorer.}\par\medskip
\medskip

On considère un repère orthonormé $\left(\text{A}~;~\vect{\imath},~\vect{\jmath},~\vect{k}\right)$ de l'espace dans lequel on place les points
\begin{center}
B(4~;~0~;~0),\quad D(0~;~4~;~0),\quad E(0~;~0~;~4) \end{center}
et les points C, F{}, G et H de sorte que le solide ABCDEFGH soit un cube.

\begin{center}
\psset{unit=0.8cm}
\begin{pspicture}(8,8.6)

\psframe(0.5,0.5)(6.1,6.1)
\psline(6.1,0.5)(7.9,2.4)(7.9,8)(6.1,6.1)
\psline(7.9,8)(2.3,8)(0.5,6.1)
\psline(0.5,6.1)(0.5,8.6)
\psline(0,8.753)(1.7,7.38)
\psline[linestyle=dashed](1.7,7.38)(8,2.321)
\psline(7.9,2.4)(8,2.321)
\psline[linestyle=dashed](0.5,0.5)(2.3,2.4)(7.9,2.4)
\psline[linestyle=dashed](2.3,2.4)(2.3,8)
\uput[dl](0.5,0.5){A} \uput[d](6.1,0.5){B} \uput[ur](7.9,2.4){C} \uput[ur](2.3,2.4){D}
\uput[l](0.5,6.1){E} \uput[r](6.1,6.1){F} \uput[ur](7.9,8){G} \uput[u](2.3,8){H} 
\uput[u](3.3,6.1){I} \uput[ul](1.4,1.45){K}
\psdots(3.3,6.1)(1.4,1.45)
\psline[linewidth=1.25pt]{->}(0.5,0.5)(1.9,0.5)\uput[d](1.2,0.5){$\vect{\imath}$}
\psline[linewidth=1.25pt]{->}(0.5,0.5)(0.95,0.975)\uput[r](0.725,0.725){$\vect{\jmath}$}
\psline[linewidth=1.25pt]{->}(0.5,0.5)(0.5,1.9)\uput[l](0.5,1.2){$\vect{k}$}
\end{pspicture}
\end{center}

\begin{enumerate}
\item Donner les coordonnées des points C, F{}, G et H.
\item On considère le point I milieu de l'arête [EF].

Montrer qu'une représentation paramétrique de la droite (IC) est donnée par: 

\[\left\{\begin{array}{l c l}
x &=&2+2t\\
y &=&\phantom{2 +~}4t\\
z &=&4 - 4t
\end{array}\right.\:\text{où} \:  t \in \R.\]

\item On désigne par $\mathcal{P}$ le plan orthogonal à la droite (IC) passant par le point G, et par J l'intersection de $\mathcal{P}$ avec (IC).
	\begin{enumerate}
		\item Démontrer qu'une équation cartésienne du plan $\mathcal{P}$ est donnée par : 
		\[x + 2y - 2z - 4 = 0.\]
		\item Justifier que J a pour coordonnées $\left(\dfrac{28}{9}~;~\dfrac{20}{9}~;~\dfrac{16}{9}\right)$.
		
Que représente J par rapport à C ? 
		\item Vérifier que le point K(0~;~2~;~0) appartient au plan $\mathcal{P}$.
		\item Justifier que (BK) est l'intersection des plans $\mathcal{P}$ et (ABC).
	\end{enumerate}
\item On rappelle que le volume d'une pyramide est donné par la formule $V = \dfrac{B \times h}{3}$, où $B$ est l'aire d'une base et $h$ la longueur de la hauteur relative à cette base.
	\begin{enumerate}
		\item Déterminer le volume de la pyramide CBKG.
		\item En déduire que l'aire du triangle BKG est égale à $12$.
		\item Justifier que la droite (BG) est incluse dans $\mathcal{P}$.
		\item On note I$'$ un point de l'arête [EF], et $P'$ le plan orthogonal à la droite (I$'$ C) passant par G.
		
Peut-on affirmer que la droite (BG) est incluse dans $P'$ ?
	\end{enumerate}
\end{enumerate}

\bigskip

\end{document}
